% =============================================================================
% projekt.tex – Selection-Panel-Projekt: Gesamtübersicht
% Einbindung: % =============================================================================
% projekt.tex – Selection-Panel-Projekt: Gesamtübersicht
% Einbindung: % =============================================================================
% projekt.tex – Selection-Panel-Projekt: Gesamtübersicht
% Einbindung: % =============================================================================
% projekt.tex – Selection-Panel-Projekt: Gesamtübersicht
% Einbindung: \input{content/projekt}
% =============================================================================

\section{Selection-Panel-Projekt -- Gesamtübersicht}
\label{sec:projekt-uebersicht}

% -----------------------------------------------------------------------------
\subsection{Begriffsrahmen und Zielsetzung}
\label{subsec:begriffsrahmen}

\textbf{Projektroot:} \code{/Users/jan/daten/IoT/embedded\_projekt}

\textbf{Zielbild:} Ein interaktives Selection Panel mit 10--100 Buttons und LEDs,
bei dem zu jedem Zeitpunkt höchstens eine Auswahl aktiv ist (One-Hot-Prinzip).
Das System besteht aus drei Ebenen: Firmware (ESP32-S3), Bridge (Raspberry Pi),
Dashboard (Browser).

\begin{table}[H]
\centering
\caption{Zentrale Begriffe}
\label{tab:projekt-begriffe}
\begin{tabularx}{\textwidth}{@{}l X@{}}
\toprule
\textbf{Begriff} & \textbf{Definition} \\
\midrule
Selection Panel & Hardware-Interface mit Tastern und LEDs für Kanalauswahl \\
One-Hot         & Invariante: genau 0 oder 1 Auswahl aktiv, niemals mehrere \\
CD4021          & 8-Bit Parallel-zu-Seriell Schieberegister (Button-Input) \\
74HC595         & 8-Bit Seriell-zu-Parallel Schieberegister (LED-Output) \\
Bridge          & Python-Server, der WebSocket $\leftrightarrow$ Serial übersetzt \\
SSOT            & Single Source of Truth -- jede Konfiguration existiert an genau einer Stelle \\
\bottomrule
\end{tabularx}
\end{table}

% -----------------------------------------------------------------------------
\subsection{Systemarchitektur -- Drei Ebenen}
\label{subsec:systemarchitektur}

Die Architektur folgt einer Dreischichten-Struktur: Browser-Dashboard kommuniziert
via WebSocket mit der Bridge, die Bridge übersetzt zwischen JSON und binärem
Serial-Protokoll zur Firmware.

\begin{figure}[H]
\centering
\begin{tikzpicture}[
  layer/.style={draw, rectangle, rounded corners=3pt, minimum width=12cm,
                minimum height=2.0cm, font=\sffamily\small},
  arrow/.style={-{Stealth[length=3mm]}, thick},
  label/.style={font=\sffamily\footnotesize, text=MutedText}
]
  % Browser Layer
  \node[layer, fill=ArduinoTeal!10] (browser) at (0,6.5) {};
  \node[font=\sffamily\bfseries] at (-4.5,7.1) {BROWSER (Dashboard)};
  \node[font=\ttfamily\footnotesize] at (0,6.8) {React + TypeScript + Vite + Tailwind};
  \node[font=\ttfamily\footnotesize] at (0,6.2) {Quellcode: apps/dashboard/src/};

  % Bridge Layer
  \node[layer, fill=RaspberryRed!10] (bridge) at (0,3.25) {};
  \node[font=\sffamily\bfseries] at (-4.5,3.85) {RASPBERRY PI (Bridge)};
  \node[font=\ttfamily\footnotesize] at (0,3.55) {Python + aiohttp};
  \node[font=\ttfamily\footnotesize] at (0,3.0) {Runtime: apps/bridge/bridge.py};
  \node[font=\ttfamily\footnotesize] at (0,2.5) {Static: apps/dashboard/dist/};

  % Firmware Layer
  \node[layer, fill=SuccessGreen!10] (firmware) at (0,0) {};
  \node[font=\sffamily\bfseries] at (-4.5,0.6) {ESP32-S3 (Firmware)};
  \node[font=\ttfamily\footnotesize] at (0,0.3) {C++ + FreeRTOS};
  \node[font=\ttfamily\footnotesize] at (0,-0.3) {Hardware: CD4021 (Buttons) + 74HC595 (LEDs)};

  % Arrows with labels
  \draw[arrow] (browser.south) -- node[right, label] {WebSocket (JSON)}
               node[left, label] {ws://rover.local:8081/ws} (bridge.north);
  \draw[arrow] (bridge.south) -- node[right, label] {USB-CDC Serial (COBS + CRC16)}
               node[left, label] {\SI{921600}{Baud}} (firmware.north);
\end{tikzpicture}
\caption{Dreischichten-Architektur des Selection Panels}
\label{fig:architektur}
\end{figure}

\paragraph{Datenfluss bei Button-Druck}

Der Ablauf vom physischen Tastendruck bis zur UI-Aktualisierung durchläuft fünf Stufen:

\begin{enumerate}
  \item CD4021 liest Button-Matrix (alle \SI{10}{\milli\second})
  \item Firmware erkennt Press-Edge, aktualisiert \code{selectedId}
  \item Firmware sendet \code{SELECTION\_STATE} via Serial (COBS/CRC)
  \item Bridge empfängt, wandelt in JSON: \code{\{"type":"selection","selected":7\}}
  \item Browser empfängt via WebSocket, React aktualisiert UI
\end{enumerate}

% -----------------------------------------------------------------------------
\subsection{Repository-Struktur}
\label{subsec:repo-struktur}

Das Repository gliedert sich in drei Hauptbereiche: Firmware, Apps (Bridge und Dashboard)
sowie Hardware-Designs und Dokumentation.

\begin{lstlisting}[style=shell, caption={Repository-Struktur}, label={lst:repo-struktur}]
embedded_projekt/
|-- apps/
|   |-- bridge/                    # Python Bridge (Pi Runtime)
|   |   |-- bridge.py             # HTTP + WebSocket + Serial
|   |   |-- docs/
|   |   |   |-- PROTOCOL.md       # Serial/WS Protokoll (SSOT)
|   |   |   |-- WORKFLOW.md       # Build/Deploy/Run
|   |   |   `-- SMOKE_TEST.md     # E2E-Tests
|   |   `-- requirements.txt
|   |
|   `-- dashboard/                 # React Frontend
|       |-- src/
|       |   |-- pages/            # DashboardPage, SettingsPage
|       |   |-- components/       # UI-Komponenten
|       |   `-- lib/              # protocol.ts, ws.ts
|       |-- dist/                 # Build-Output (wird von Bridge served)
|       |-- docs/
|       |   |-- DEVELOPER.md      # Frontend-Entwicklung
|       |   `-- USER_GUIDE.md     # Benutzerhandbuch
|       |-- vite.config.ts
|       `-- package.json
|
|-- firmware/                      # ESP32-S3 Firmware (C++)
|   |-- include/hw/               # SSOT: Pins, Counts, SPI-Parameter
|   |-- src/
|   |   |-- hal/                  # SPI-Bus Zugriff
|   |   |-- drivers/              # CD4021, HC595
|   |   |-- logic/                # Debounce, Selection, Pipeline
|   |   |-- app/                  # Loop/Task Orchestrierung
|   |   `-- protocol/             # COBS/CRC Frames
|   |-- test/                     # Unit + E2E Tests
|   |-- docs/                     # Firmware-Dokumentation
|   `-- platformio.ini
|
|-- hardware/
|   `-- kicad/                    # Schaltplaene, PCB
|
|-- doc/
|   |-- md/                       # Markdown-Quellen (SSOT)
|   `-- latex/                    # LaTeX/PDF-Output
|
|-- assets/
|   |-- pdf/                      # Generierte PDFs
|   `-- video/                    # Medien
|
`-- scripts/                       # Build-/Deploy-Helfer
\end{lstlisting}

% -----------------------------------------------------------------------------
\subsection{SSOT -- Single Source of Truth}
\label{subsec:ssot}

\textbf{Regel:} Jede Konfiguration existiert an genau einer Stelle. Änderungen
propagieren automatisch.

\begin{table}[H]
\centering
\caption{SSOT-Dateien nach Aspekt}
\label{tab:ssot}
\begin{tabularx}{\textwidth}{@{}l X@{}}
\toprule
\textbf{Aspekt} & \textbf{SSOT-Datei} \\
\midrule
Pins                & \code{firmware/include/hw/board\_pins.h} \\
Counts/Masks        & \code{firmware/include/hw/io\_counts.h} \\
SPI-Parameter       & \code{firmware/include/hw/spi\_params.h} \\
Serial-Protokoll    & \code{apps/bridge/docs/PROTOCOL.md} \\
WebSocket-Protokoll & \code{apps/bridge/docs/PROTOCOL.md} \\
\bottomrule
\end{tabularx}
\end{table}

\begin{table}[H]
\centering
\caption{Generierte Artefakte}
\label{tab:artefakte}
\begin{tabularx}{\textwidth}{@{}l l X@{}}
\toprule
\textbf{Artefakt} & \textbf{Quelle} & \textbf{Output} \\
\midrule
Frontend & \code{apps/dashboard/src/} & \code{apps/dashboard/dist/} \\
PDFs     & \code{doc/latex/}          & \code{assets/pdf/} \\
\bottomrule
\end{tabularx}
\end{table}

% -----------------------------------------------------------------------------
\subsection{Firmware -- Build, Flash, Test}
\label{subsec:firmware}

\paragraph{Flashen}

\begin{lstlisting}[style=shell, caption={Firmware flashen}, label={lst:flash}]
cd /Users/jan/daten/IoT/embedded_projekt

# FreeRTOS-Variante (empfohlen)
pio run -d firmware -e seeed_xiao_esp32s3_freertos -t upload

# Arduino-Loop-Variante
pio run -d firmware -e seeed_xiao_esp32s3_loop -t upload
\end{lstlisting}

\paragraph{Monitor}

\begin{lstlisting}[style=shell]
pio device monitor -d firmware -e seeed_xiao_esp32s3_freertos
\end{lstlisting}

\paragraph{Tests}

\begin{lstlisting}[style=shell, caption={Unit- und E2E-Tests}, label={lst:tests}]
# Unit-Tests (Logik: Debounce, Selection, Pipeline)
pio test -d firmware -e seeed_xiao_esp32s3_test -v

# E2E-Hardwaretest (SPI, Pinning, Timing)
pio test -d firmware -e seeed_xiao_esp32s3_e2e -f test_e2e_hw -v
\end{lstlisting}

\begin{table}[H]
\centering
\caption{Build-Varianten}
\label{tab:build-varianten}
\begin{tabularx}{\textwidth}{@{}l X@{}}
\toprule
\textbf{Environment} & \textbf{Beschreibung} \\
\midrule
\code{seeed\_xiao\_esp32s3\_freertos} & FreeRTOS Task (\code{APP\_USE\_FREERTOS=1}) \\
\code{seeed\_xiao\_esp32s3\_loop}     & Arduino Loop (\code{APP\_USE\_FREERTOS=0}) \\
\code{seeed\_xiao\_esp32s3\_test}     & Unit-Tests (nur \code{logic/}) \\
\code{seeed\_xiao\_esp32s3\_e2e}      & E2E-Hardware-Test \\
\bottomrule
\end{tabularx}
\end{table}

% -----------------------------------------------------------------------------
\subsection{Dashboard -- Entwicklung und Build}
\label{subsec:dashboard}

\paragraph{DEV-Server (Mac)}

\begin{lstlisting}[style=shell]
cd /Users/jan/daten/IoT/embedded_projekt/apps/dashboard
npm ci
npm run dev
# -> http://localhost:5173/
# -> WebSocket via Proxy zu rover.local:8081
\end{lstlisting}

\paragraph{Production-Build}

\begin{lstlisting}[style=shell]
npm run lint
npm run typecheck
npm run build
# -> Output: dist/
\end{lstlisting}

\begin{table}[H]
\centering
\caption{npm Scripts}
\label{tab:npm-scripts}
\begin{tabularx}{\textwidth}{@{}l X@{}}
\toprule
\textbf{Script} & \textbf{Zweck} \\
\midrule
\code{npm run dev}       & DEV-Server mit HMR \\
\code{npm run lint}      & ESLint-Prüfung \\
\code{npm run typecheck} & TypeScript-Prüfung \\
\code{npm run build}     & Production-Build nach \code{dist/} \\
\code{npm run format}    & Prettier-Formatierung \\
\bottomrule
\end{tabularx}
\end{table}

% -----------------------------------------------------------------------------
\subsection{Bridge -- Raspberry Pi}
\label{subsec:bridge}

\paragraph{Manueller Start}

\begin{lstlisting}[style=shell, caption={Bridge manuell starten}, label={lst:bridge-start}]
ssh pi@rover
source ~/venvs/selection-bridge/bin/activate

python /home/pi/embedded_projekt/apps/bridge/bridge.py \
  --port auto \
  --vid 0x303A --pid 0x1001 \
  --baud 921600 \
  --http 8081 \
  --web-root /home/pi/embedded_projekt/apps/dashboard/dist \
  --wait-port --wait-timeout 0 --wait-poll 0.5 \
  --reconnect-delay 2 \
  --ack-timeout 0.6 --retries 4
\end{lstlisting}

\paragraph{systemd-Service}

\begin{lstlisting}[style=shell, caption={systemd-Befehle}, label={lst:systemd}]
# Status pruefen
sudo systemctl status embedded-bridge.service --no-pager -l

# Neu starten
sudo systemctl restart embedded-bridge.service

# Logs verfolgen
sudo journalctl -u embedded-bridge.service -f
\end{lstlisting}

\begin{table}[H]
\centering
\caption{URLs im PROD-Modus}
\label{tab:urls}
\begin{tabularx}{\textwidth}{@{}l X@{}}
\toprule
\textbf{Zweck} & \textbf{URL} \\
\midrule
Dashboard & \url{http://rover.local:8081/} \\
WebSocket & \url{ws://rover.local:8081/ws} \\
\bottomrule
\end{tabularx}
\end{table}

% -----------------------------------------------------------------------------
\subsection{Deployment-Workflow}
\label{subsec:deployment}

\paragraph{Mac → Pi synchronisieren}

\begin{lstlisting}[style=shell, caption={rsync-Synchronisation}, label={lst:rsync}]
# Zielpfad auf dem Pi anlegen (einmalig)
ssh pi@rover 'mkdir -p /home/pi/embedded_projekt/apps/bridge \
                       /home/pi/embedded_projekt/apps/dashboard'

# Sync (ohne venv, node_modules, dist)
rsync -avh --delete \
  --exclude '.git/' \
  --exclude '.venv/' \
  --exclude '__pycache__/' \
  --exclude '*.pyc' \
  --exclude '.DS_Store' \
  --exclude 'node_modules/' \
  --exclude 'apps/dashboard/dist/' \
  ./ \
  pi@rover:/home/pi/embedded_projekt/
\end{lstlisting}

\paragraph{Dashboard dist sync}

\begin{lstlisting}[style=shell, caption={Dashboard-Build und Sync}, label={lst:dist-sync}]
# Lokal bauen
cd apps/dashboard
npm ci
npm run build
cd -

# Dashboard dist sync (prod)
rsync -avh --delete \
  --exclude '.DS_Store' \
  ./apps/dashboard/dist/ \
  pi@rover:/home/pi/embedded_projekt/apps/dashboard/dist/
\end{lstlisting}

\paragraph{Service neu starten}

\begin{lstlisting}[style=shell]
ssh pi@rover "sudo systemctl daemon-reload && \
              sudo systemctl restart embedded-bridge.service"
\end{lstlisting}

% -----------------------------------------------------------------------------
\subsection{Dokumentationsstruktur}
\label{subsec:dokumentation}

\paragraph{Firmware-Dokumentation}

Die Firmware-Dokumentation liegt unter \code{firmware/docs/}:

\begin{table}[H]
\centering
\caption{Firmware-Dokumentation}
\label{tab:firmware-docs}
\begin{tabularx}{\textwidth}{@{}l X@{}}
\toprule
\textbf{Datei} & \textbf{Inhalt} \\
\midrule
\code{overview.md}     & Projektübersicht, Bit-Mapping (MSB-first) \\
\code{wiring.md}       & Hardware-Verdrahtung (Pins, ICs) \\
\code{architecture.md} & Schichtenmodell (HAL → Driver → Logic → App) \\
\bottomrule
\end{tabularx}
\end{table}

\paragraph{Bridge-Dokumentation}

Die Bridge-Dokumentation liegt unter \code{apps/bridge/docs/}:

\begin{table}[H]
\centering
\caption{Bridge-Dokumentation}
\label{tab:bridge-docs}
\begin{tabularx}{\textwidth}{@{}l X@{}}
\toprule
\textbf{Datei} & \textbf{Inhalt} \\
\midrule
\code{PROTOCOL.md}   & Serial + WebSocket Protokoll (SSOT) \\
\code{WORKFLOW.md}   & Build/Deploy/Run \\
\code{SMOKE\_TEST.md} & E2E-Tests \\
\bottomrule
\end{tabularx}
\end{table}

\paragraph{Dashboard-Dokumentation}

Die Dashboard-Dokumentation liegt unter \code{apps/dashboard/docs/}:

\begin{table}[H]
\centering
\caption{Dashboard-Dokumentation}
\label{tab:dashboard-docs}
\begin{tabularx}{\textwidth}{@{}l X@{}}
\toprule
\textbf{Datei} & \textbf{Inhalt} \\
\midrule
\code{DEVELOPER.md}  & Frontend-Entwicklung \\
\code{USER\_GUIDE.md} & Benutzerhandbuch \\
\bottomrule
\end{tabularx}
\end{table}

% -----------------------------------------------------------------------------
\subsection{Typische Arbeitsabläufe}
\label{subsec:workflows}

\begin{infobox}[Firmware-Logik ändern]
\begin{enumerate}
  \item Code ändern in \code{firmware/src/logic/}
  \item Unit-Tests: \code{pio test -d firmware -e seeed\_xiao\_esp32s3\_test -v}
  \item Bei Erfolg: Flashen mit \code{pio run ... -t upload}
  \item Monitor: \code{pio device monitor -d firmware}
\end{enumerate}
\end{infobox}

\begin{infobox}[Dashboard-UI ändern]
\begin{enumerate}
  \item DEV-Server starten: \code{cd apps/dashboard \&\& npm run dev}
  \item Ändern -- Browser aktualisiert automatisch (HMR)
  \item Lint + Build: \code{npm run lint \&\& npm run build}
  \item Deploy (siehe \cref{subsec:deployment})
\end{enumerate}
\end{infobox}

\begin{warnbox}[Hardware-Verdrahtung ändern]
Bei Änderungen an der Hardware-Verdrahtung:
\begin{enumerate}
  \item SSOT aktualisieren: \code{firmware/include/hw/board\_pins.h}
  \item Dokumentation aktualisieren: \code{firmware/docs/wiring.md}
  \item E2E-Test: \code{pio test ... -e seeed\_xiao\_esp32s3\_e2e -f test\_e2e\_hw -v}
\end{enumerate}
\end{warnbox}

% -----------------------------------------------------------------------------
\subsection{Roadmap}
\label{subsec:roadmap}

\begin{table}[H]
\centering
\caption{Projekt-Roadmap}
\label{tab:roadmap}
\begin{tabularx}{\textwidth}{@{}l c X@{}}
\toprule
\textbf{Milestone} & \textbf{Status} & \textbf{Beschreibung} \\
\midrule
10× Buttons/LEDs   & ✓  & Grundfunktion verifiziert \\
Serial-Protokoll   & ✓  & COBS/CRC, ACK/Retry \\
React-Dashboard    & ✓  & DEV/PROD-Split \\
100× Skalierung    & -- & Counts/Masks + Kaskadierung + Timing \\
Audio-Integration  & -- & Multimedia bei Selection-Change \\
\bottomrule
\end{tabularx}
\end{table}

% -----------------------------------------------------------------------------
\subsection{Siehe auch}
\label{subsec:projekt-siehe-auch}

Weiterführende Dokumentation:

\begin{itemize}
  \item \code{firmware/docs/overview.md} -- Firmware-Einstieg
  \item \code{apps/dashboard/docs/DEVELOPER.md} -- Frontend-Einstieg
  \item \code{apps/bridge/docs/WORKFLOW.md} -- Pi-Deployment
  \item \code{apps/bridge/docs/PROTOCOL.md} -- Protokoll-Spezifikation
\end{itemize}

% =============================================================================

\section{Selection-Panel-Projekt -- Gesamtübersicht}
\label{sec:projekt-uebersicht}

% -----------------------------------------------------------------------------
\subsection{Begriffsrahmen und Zielsetzung}
\label{subsec:begriffsrahmen}

\textbf{Projektroot:} \code{/Users/jan/daten/IoT/embedded\_projekt}

\textbf{Zielbild:} Ein interaktives Selection Panel mit 10--100 Buttons und LEDs,
bei dem zu jedem Zeitpunkt höchstens eine Auswahl aktiv ist (One-Hot-Prinzip).
Das System besteht aus drei Ebenen: Firmware (ESP32-S3), Bridge (Raspberry Pi),
Dashboard (Browser).

\begin{table}[H]
\centering
\caption{Zentrale Begriffe}
\label{tab:projekt-begriffe}
\begin{tabularx}{\textwidth}{@{}l X@{}}
\toprule
\textbf{Begriff} & \textbf{Definition} \\
\midrule
Selection Panel & Hardware-Interface mit Tastern und LEDs für Kanalauswahl \\
One-Hot         & Invariante: genau 0 oder 1 Auswahl aktiv, niemals mehrere \\
CD4021          & 8-Bit Parallel-zu-Seriell Schieberegister (Button-Input) \\
74HC595         & 8-Bit Seriell-zu-Parallel Schieberegister (LED-Output) \\
Bridge          & Python-Server, der WebSocket $\leftrightarrow$ Serial übersetzt \\
SSOT            & Single Source of Truth -- jede Konfiguration existiert an genau einer Stelle \\
\bottomrule
\end{tabularx}
\end{table}

% -----------------------------------------------------------------------------
\subsection{Systemarchitektur -- Drei Ebenen}
\label{subsec:systemarchitektur}

Die Architektur folgt einer Dreischichten-Struktur: Browser-Dashboard kommuniziert
via WebSocket mit der Bridge, die Bridge übersetzt zwischen JSON und binärem
Serial-Protokoll zur Firmware.

\begin{figure}[H]
\centering
\begin{tikzpicture}[
  layer/.style={draw, rectangle, rounded corners=3pt, minimum width=12cm,
                minimum height=2.0cm, font=\sffamily\small},
  arrow/.style={-{Stealth[length=3mm]}, thick},
  label/.style={font=\sffamily\footnotesize, text=MutedText}
]
  % Browser Layer
  \node[layer, fill=ArduinoTeal!10] (browser) at (0,6.5) {};
  \node[font=\sffamily\bfseries] at (-4.5,7.1) {BROWSER (Dashboard)};
  \node[font=\ttfamily\footnotesize] at (0,6.8) {React + TypeScript + Vite + Tailwind};
  \node[font=\ttfamily\footnotesize] at (0,6.2) {Quellcode: apps/dashboard/src/};

  % Bridge Layer
  \node[layer, fill=RaspberryRed!10] (bridge) at (0,3.25) {};
  \node[font=\sffamily\bfseries] at (-4.5,3.85) {RASPBERRY PI (Bridge)};
  \node[font=\ttfamily\footnotesize] at (0,3.55) {Python + aiohttp};
  \node[font=\ttfamily\footnotesize] at (0,3.0) {Runtime: apps/bridge/bridge.py};
  \node[font=\ttfamily\footnotesize] at (0,2.5) {Static: apps/dashboard/dist/};

  % Firmware Layer
  \node[layer, fill=SuccessGreen!10] (firmware) at (0,0) {};
  \node[font=\sffamily\bfseries] at (-4.5,0.6) {ESP32-S3 (Firmware)};
  \node[font=\ttfamily\footnotesize] at (0,0.3) {C++ + FreeRTOS};
  \node[font=\ttfamily\footnotesize] at (0,-0.3) {Hardware: CD4021 (Buttons) + 74HC595 (LEDs)};

  % Arrows with labels
  \draw[arrow] (browser.south) -- node[right, label] {WebSocket (JSON)}
               node[left, label] {ws://rover.local:8081/ws} (bridge.north);
  \draw[arrow] (bridge.south) -- node[right, label] {USB-CDC Serial (COBS + CRC16)}
               node[left, label] {\SI{921600}{Baud}} (firmware.north);
\end{tikzpicture}
\caption{Dreischichten-Architektur des Selection Panels}
\label{fig:architektur}
\end{figure}

\paragraph{Datenfluss bei Button-Druck}

Der Ablauf vom physischen Tastendruck bis zur UI-Aktualisierung durchläuft fünf Stufen:

\begin{enumerate}
  \item CD4021 liest Button-Matrix (alle \SI{10}{\milli\second})
  \item Firmware erkennt Press-Edge, aktualisiert \code{selectedId}
  \item Firmware sendet \code{SELECTION\_STATE} via Serial (COBS/CRC)
  \item Bridge empfängt, wandelt in JSON: \code{\{"type":"selection","selected":7\}}
  \item Browser empfängt via WebSocket, React aktualisiert UI
\end{enumerate}

% -----------------------------------------------------------------------------
\subsection{Repository-Struktur}
\label{subsec:repo-struktur}

Das Repository gliedert sich in drei Hauptbereiche: Firmware, Apps (Bridge und Dashboard)
sowie Hardware-Designs und Dokumentation.

\begin{lstlisting}[style=shell, caption={Repository-Struktur}, label={lst:repo-struktur}]
embedded_projekt/
|-- apps/
|   |-- bridge/                    # Python Bridge (Pi Runtime)
|   |   |-- bridge.py             # HTTP + WebSocket + Serial
|   |   |-- docs/
|   |   |   |-- PROTOCOL.md       # Serial/WS Protokoll (SSOT)
|   |   |   |-- WORKFLOW.md       # Build/Deploy/Run
|   |   |   `-- SMOKE_TEST.md     # E2E-Tests
|   |   `-- requirements.txt
|   |
|   `-- dashboard/                 # React Frontend
|       |-- src/
|       |   |-- pages/            # DashboardPage, SettingsPage
|       |   |-- components/       # UI-Komponenten
|       |   `-- lib/              # protocol.ts, ws.ts
|       |-- dist/                 # Build-Output (wird von Bridge served)
|       |-- docs/
|       |   |-- DEVELOPER.md      # Frontend-Entwicklung
|       |   `-- USER_GUIDE.md     # Benutzerhandbuch
|       |-- vite.config.ts
|       `-- package.json
|
|-- firmware/                      # ESP32-S3 Firmware (C++)
|   |-- include/hw/               # SSOT: Pins, Counts, SPI-Parameter
|   |-- src/
|   |   |-- hal/                  # SPI-Bus Zugriff
|   |   |-- drivers/              # CD4021, HC595
|   |   |-- logic/                # Debounce, Selection, Pipeline
|   |   |-- app/                  # Loop/Task Orchestrierung
|   |   `-- protocol/             # COBS/CRC Frames
|   |-- test/                     # Unit + E2E Tests
|   |-- docs/                     # Firmware-Dokumentation
|   `-- platformio.ini
|
|-- hardware/
|   `-- kicad/                    # Schaltplaene, PCB
|
|-- doc/
|   |-- md/                       # Markdown-Quellen (SSOT)
|   `-- latex/                    # LaTeX/PDF-Output
|
|-- assets/
|   |-- pdf/                      # Generierte PDFs
|   `-- video/                    # Medien
|
`-- scripts/                       # Build-/Deploy-Helfer
\end{lstlisting}

% -----------------------------------------------------------------------------
\subsection{SSOT -- Single Source of Truth}
\label{subsec:ssot}

\textbf{Regel:} Jede Konfiguration existiert an genau einer Stelle. Änderungen
propagieren automatisch.

\begin{table}[H]
\centering
\caption{SSOT-Dateien nach Aspekt}
\label{tab:ssot}
\begin{tabularx}{\textwidth}{@{}l X@{}}
\toprule
\textbf{Aspekt} & \textbf{SSOT-Datei} \\
\midrule
Pins                & \code{firmware/include/hw/board\_pins.h} \\
Counts/Masks        & \code{firmware/include/hw/io\_counts.h} \\
SPI-Parameter       & \code{firmware/include/hw/spi\_params.h} \\
Serial-Protokoll    & \code{apps/bridge/docs/PROTOCOL.md} \\
WebSocket-Protokoll & \code{apps/bridge/docs/PROTOCOL.md} \\
\bottomrule
\end{tabularx}
\end{table}

\begin{table}[H]
\centering
\caption{Generierte Artefakte}
\label{tab:artefakte}
\begin{tabularx}{\textwidth}{@{}l l X@{}}
\toprule
\textbf{Artefakt} & \textbf{Quelle} & \textbf{Output} \\
\midrule
Frontend & \code{apps/dashboard/src/} & \code{apps/dashboard/dist/} \\
PDFs     & \code{doc/latex/}          & \code{assets/pdf/} \\
\bottomrule
\end{tabularx}
\end{table}

% -----------------------------------------------------------------------------
\subsection{Firmware -- Build, Flash, Test}
\label{subsec:firmware}

\paragraph{Flashen}

\begin{lstlisting}[style=shell, caption={Firmware flashen}, label={lst:flash}]
cd /Users/jan/daten/IoT/embedded_projekt

# FreeRTOS-Variante (empfohlen)
pio run -d firmware -e seeed_xiao_esp32s3_freertos -t upload

# Arduino-Loop-Variante
pio run -d firmware -e seeed_xiao_esp32s3_loop -t upload
\end{lstlisting}

\paragraph{Monitor}

\begin{lstlisting}[style=shell]
pio device monitor -d firmware -e seeed_xiao_esp32s3_freertos
\end{lstlisting}

\paragraph{Tests}

\begin{lstlisting}[style=shell, caption={Unit- und E2E-Tests}, label={lst:tests}]
# Unit-Tests (Logik: Debounce, Selection, Pipeline)
pio test -d firmware -e seeed_xiao_esp32s3_test -v

# E2E-Hardwaretest (SPI, Pinning, Timing)
pio test -d firmware -e seeed_xiao_esp32s3_e2e -f test_e2e_hw -v
\end{lstlisting}

\begin{table}[H]
\centering
\caption{Build-Varianten}
\label{tab:build-varianten}
\begin{tabularx}{\textwidth}{@{}l X@{}}
\toprule
\textbf{Environment} & \textbf{Beschreibung} \\
\midrule
\code{seeed\_xiao\_esp32s3\_freertos} & FreeRTOS Task (\code{APP\_USE\_FREERTOS=1}) \\
\code{seeed\_xiao\_esp32s3\_loop}     & Arduino Loop (\code{APP\_USE\_FREERTOS=0}) \\
\code{seeed\_xiao\_esp32s3\_test}     & Unit-Tests (nur \code{logic/}) \\
\code{seeed\_xiao\_esp32s3\_e2e}      & E2E-Hardware-Test \\
\bottomrule
\end{tabularx}
\end{table}

% -----------------------------------------------------------------------------
\subsection{Dashboard -- Entwicklung und Build}
\label{subsec:dashboard}

\paragraph{DEV-Server (Mac)}

\begin{lstlisting}[style=shell]
cd /Users/jan/daten/IoT/embedded_projekt/apps/dashboard
npm ci
npm run dev
# -> http://localhost:5173/
# -> WebSocket via Proxy zu rover.local:8081
\end{lstlisting}

\paragraph{Production-Build}

\begin{lstlisting}[style=shell]
npm run lint
npm run typecheck
npm run build
# -> Output: dist/
\end{lstlisting}

\begin{table}[H]
\centering
\caption{npm Scripts}
\label{tab:npm-scripts}
\begin{tabularx}{\textwidth}{@{}l X@{}}
\toprule
\textbf{Script} & \textbf{Zweck} \\
\midrule
\code{npm run dev}       & DEV-Server mit HMR \\
\code{npm run lint}      & ESLint-Prüfung \\
\code{npm run typecheck} & TypeScript-Prüfung \\
\code{npm run build}     & Production-Build nach \code{dist/} \\
\code{npm run format}    & Prettier-Formatierung \\
\bottomrule
\end{tabularx}
\end{table}

% -----------------------------------------------------------------------------
\subsection{Bridge -- Raspberry Pi}
\label{subsec:bridge}

\paragraph{Manueller Start}

\begin{lstlisting}[style=shell, caption={Bridge manuell starten}, label={lst:bridge-start}]
ssh pi@rover
source ~/venvs/selection-bridge/bin/activate

python /home/pi/embedded_projekt/apps/bridge/bridge.py \
  --port auto \
  --vid 0x303A --pid 0x1001 \
  --baud 921600 \
  --http 8081 \
  --web-root /home/pi/embedded_projekt/apps/dashboard/dist \
  --wait-port --wait-timeout 0 --wait-poll 0.5 \
  --reconnect-delay 2 \
  --ack-timeout 0.6 --retries 4
\end{lstlisting}

\paragraph{systemd-Service}

\begin{lstlisting}[style=shell, caption={systemd-Befehle}, label={lst:systemd}]
# Status pruefen
sudo systemctl status embedded-bridge.service --no-pager -l

# Neu starten
sudo systemctl restart embedded-bridge.service

# Logs verfolgen
sudo journalctl -u embedded-bridge.service -f
\end{lstlisting}

\begin{table}[H]
\centering
\caption{URLs im PROD-Modus}
\label{tab:urls}
\begin{tabularx}{\textwidth}{@{}l X@{}}
\toprule
\textbf{Zweck} & \textbf{URL} \\
\midrule
Dashboard & \url{http://rover.local:8081/} \\
WebSocket & \url{ws://rover.local:8081/ws} \\
\bottomrule
\end{tabularx}
\end{table}

% -----------------------------------------------------------------------------
\subsection{Deployment-Workflow}
\label{subsec:deployment}

\paragraph{Mac → Pi synchronisieren}

\begin{lstlisting}[style=shell, caption={rsync-Synchronisation}, label={lst:rsync}]
# Zielpfad auf dem Pi anlegen (einmalig)
ssh pi@rover 'mkdir -p /home/pi/embedded_projekt/apps/bridge \
                       /home/pi/embedded_projekt/apps/dashboard'

# Sync (ohne venv, node_modules, dist)
rsync -avh --delete \
  --exclude '.git/' \
  --exclude '.venv/' \
  --exclude '__pycache__/' \
  --exclude '*.pyc' \
  --exclude '.DS_Store' \
  --exclude 'node_modules/' \
  --exclude 'apps/dashboard/dist/' \
  ./ \
  pi@rover:/home/pi/embedded_projekt/
\end{lstlisting}

\paragraph{Dashboard dist sync}

\begin{lstlisting}[style=shell, caption={Dashboard-Build und Sync}, label={lst:dist-sync}]
# Lokal bauen
cd apps/dashboard
npm ci
npm run build
cd -

# Dashboard dist sync (prod)
rsync -avh --delete \
  --exclude '.DS_Store' \
  ./apps/dashboard/dist/ \
  pi@rover:/home/pi/embedded_projekt/apps/dashboard/dist/
\end{lstlisting}

\paragraph{Service neu starten}

\begin{lstlisting}[style=shell]
ssh pi@rover "sudo systemctl daemon-reload && \
              sudo systemctl restart embedded-bridge.service"
\end{lstlisting}

% -----------------------------------------------------------------------------
\subsection{Dokumentationsstruktur}
\label{subsec:dokumentation}

\paragraph{Firmware-Dokumentation}

Die Firmware-Dokumentation liegt unter \code{firmware/docs/}:

\begin{table}[H]
\centering
\caption{Firmware-Dokumentation}
\label{tab:firmware-docs}
\begin{tabularx}{\textwidth}{@{}l X@{}}
\toprule
\textbf{Datei} & \textbf{Inhalt} \\
\midrule
\code{overview.md}     & Projektübersicht, Bit-Mapping (MSB-first) \\
\code{wiring.md}       & Hardware-Verdrahtung (Pins, ICs) \\
\code{architecture.md} & Schichtenmodell (HAL → Driver → Logic → App) \\
\bottomrule
\end{tabularx}
\end{table}

\paragraph{Bridge-Dokumentation}

Die Bridge-Dokumentation liegt unter \code{apps/bridge/docs/}:

\begin{table}[H]
\centering
\caption{Bridge-Dokumentation}
\label{tab:bridge-docs}
\begin{tabularx}{\textwidth}{@{}l X@{}}
\toprule
\textbf{Datei} & \textbf{Inhalt} \\
\midrule
\code{PROTOCOL.md}   & Serial + WebSocket Protokoll (SSOT) \\
\code{WORKFLOW.md}   & Build/Deploy/Run \\
\code{SMOKE\_TEST.md} & E2E-Tests \\
\bottomrule
\end{tabularx}
\end{table}

\paragraph{Dashboard-Dokumentation}

Die Dashboard-Dokumentation liegt unter \code{apps/dashboard/docs/}:

\begin{table}[H]
\centering
\caption{Dashboard-Dokumentation}
\label{tab:dashboard-docs}
\begin{tabularx}{\textwidth}{@{}l X@{}}
\toprule
\textbf{Datei} & \textbf{Inhalt} \\
\midrule
\code{DEVELOPER.md}  & Frontend-Entwicklung \\
\code{USER\_GUIDE.md} & Benutzerhandbuch \\
\bottomrule
\end{tabularx}
\end{table}

% -----------------------------------------------------------------------------
\subsection{Typische Arbeitsabläufe}
\label{subsec:workflows}

\begin{infobox}[Firmware-Logik ändern]
\begin{enumerate}
  \item Code ändern in \code{firmware/src/logic/}
  \item Unit-Tests: \code{pio test -d firmware -e seeed\_xiao\_esp32s3\_test -v}
  \item Bei Erfolg: Flashen mit \code{pio run ... -t upload}
  \item Monitor: \code{pio device monitor -d firmware}
\end{enumerate}
\end{infobox}

\begin{infobox}[Dashboard-UI ändern]
\begin{enumerate}
  \item DEV-Server starten: \code{cd apps/dashboard \&\& npm run dev}
  \item Ändern -- Browser aktualisiert automatisch (HMR)
  \item Lint + Build: \code{npm run lint \&\& npm run build}
  \item Deploy (siehe \cref{subsec:deployment})
\end{enumerate}
\end{infobox}

\begin{warnbox}[Hardware-Verdrahtung ändern]
Bei Änderungen an der Hardware-Verdrahtung:
\begin{enumerate}
  \item SSOT aktualisieren: \code{firmware/include/hw/board\_pins.h}
  \item Dokumentation aktualisieren: \code{firmware/docs/wiring.md}
  \item E2E-Test: \code{pio test ... -e seeed\_xiao\_esp32s3\_e2e -f test\_e2e\_hw -v}
\end{enumerate}
\end{warnbox}

% -----------------------------------------------------------------------------
\subsection{Roadmap}
\label{subsec:roadmap}

\begin{table}[H]
\centering
\caption{Projekt-Roadmap}
\label{tab:roadmap}
\begin{tabularx}{\textwidth}{@{}l c X@{}}
\toprule
\textbf{Milestone} & \textbf{Status} & \textbf{Beschreibung} \\
\midrule
10× Buttons/LEDs   & ✓  & Grundfunktion verifiziert \\
Serial-Protokoll   & ✓  & COBS/CRC, ACK/Retry \\
React-Dashboard    & ✓  & DEV/PROD-Split \\
100× Skalierung    & -- & Counts/Masks + Kaskadierung + Timing \\
Audio-Integration  & -- & Multimedia bei Selection-Change \\
\bottomrule
\end{tabularx}
\end{table}

% -----------------------------------------------------------------------------
\subsection{Siehe auch}
\label{subsec:projekt-siehe-auch}

Weiterführende Dokumentation:

\begin{itemize}
  \item \code{firmware/docs/overview.md} -- Firmware-Einstieg
  \item \code{apps/dashboard/docs/DEVELOPER.md} -- Frontend-Einstieg
  \item \code{apps/bridge/docs/WORKFLOW.md} -- Pi-Deployment
  \item \code{apps/bridge/docs/PROTOCOL.md} -- Protokoll-Spezifikation
\end{itemize}

% =============================================================================

\section{Selection-Panel-Projekt -- Gesamtübersicht}
\label{sec:projekt-uebersicht}

% -----------------------------------------------------------------------------
\subsection{Begriffsrahmen und Zielsetzung}
\label{subsec:begriffsrahmen}

\textbf{Projektroot:} \code{/Users/jan/daten/IoT/embedded\_projekt}

\textbf{Zielbild:} Ein interaktives Selection Panel mit 10--100 Buttons und LEDs,
bei dem zu jedem Zeitpunkt höchstens eine Auswahl aktiv ist (One-Hot-Prinzip).
Das System besteht aus drei Ebenen: Firmware (ESP32-S3), Bridge (Raspberry Pi),
Dashboard (Browser).

\begin{table}[H]
\centering
\caption{Zentrale Begriffe}
\label{tab:projekt-begriffe}
\begin{tabularx}{\textwidth}{@{}l X@{}}
\toprule
\textbf{Begriff} & \textbf{Definition} \\
\midrule
Selection Panel & Hardware-Interface mit Tastern und LEDs für Kanalauswahl \\
One-Hot         & Invariante: genau 0 oder 1 Auswahl aktiv, niemals mehrere \\
CD4021          & 8-Bit Parallel-zu-Seriell Schieberegister (Button-Input) \\
74HC595         & 8-Bit Seriell-zu-Parallel Schieberegister (LED-Output) \\
Bridge          & Python-Server, der WebSocket $\leftrightarrow$ Serial übersetzt \\
SSOT            & Single Source of Truth -- jede Konfiguration existiert an genau einer Stelle \\
\bottomrule
\end{tabularx}
\end{table}

% -----------------------------------------------------------------------------
\subsection{Systemarchitektur -- Drei Ebenen}
\label{subsec:systemarchitektur}

Die Architektur folgt einer Dreischichten-Struktur: Browser-Dashboard kommuniziert
via WebSocket mit der Bridge, die Bridge übersetzt zwischen JSON und binärem
Serial-Protokoll zur Firmware.

\begin{figure}[H]
\centering
\begin{tikzpicture}[
  layer/.style={draw, rectangle, rounded corners=3pt, minimum width=12cm,
                minimum height=2.0cm, font=\sffamily\small},
  arrow/.style={-{Stealth[length=3mm]}, thick},
  label/.style={font=\sffamily\footnotesize, text=MutedText}
]
  % Browser Layer
  \node[layer, fill=ArduinoTeal!10] (browser) at (0,6.5) {};
  \node[font=\sffamily\bfseries] at (-4.5,7.1) {BROWSER (Dashboard)};
  \node[font=\ttfamily\footnotesize] at (0,6.8) {React + TypeScript + Vite + Tailwind};
  \node[font=\ttfamily\footnotesize] at (0,6.2) {Quellcode: apps/dashboard/src/};

  % Bridge Layer
  \node[layer, fill=RaspberryRed!10] (bridge) at (0,3.25) {};
  \node[font=\sffamily\bfseries] at (-4.5,3.85) {RASPBERRY PI (Bridge)};
  \node[font=\ttfamily\footnotesize] at (0,3.55) {Python + aiohttp};
  \node[font=\ttfamily\footnotesize] at (0,3.0) {Runtime: apps/bridge/bridge.py};
  \node[font=\ttfamily\footnotesize] at (0,2.5) {Static: apps/dashboard/dist/};

  % Firmware Layer
  \node[layer, fill=SuccessGreen!10] (firmware) at (0,0) {};
  \node[font=\sffamily\bfseries] at (-4.5,0.6) {ESP32-S3 (Firmware)};
  \node[font=\ttfamily\footnotesize] at (0,0.3) {C++ + FreeRTOS};
  \node[font=\ttfamily\footnotesize] at (0,-0.3) {Hardware: CD4021 (Buttons) + 74HC595 (LEDs)};

  % Arrows with labels
  \draw[arrow] (browser.south) -- node[right, label] {WebSocket (JSON)}
               node[left, label] {ws://rover.local:8081/ws} (bridge.north);
  \draw[arrow] (bridge.south) -- node[right, label] {USB-CDC Serial (COBS + CRC16)}
               node[left, label] {\SI{921600}{Baud}} (firmware.north);
\end{tikzpicture}
\caption{Dreischichten-Architektur des Selection Panels}
\label{fig:architektur}
\end{figure}

\paragraph{Datenfluss bei Button-Druck}

Der Ablauf vom physischen Tastendruck bis zur UI-Aktualisierung durchläuft fünf Stufen:

\begin{enumerate}
  \item CD4021 liest Button-Matrix (alle \SI{10}{\milli\second})
  \item Firmware erkennt Press-Edge, aktualisiert \code{selectedId}
  \item Firmware sendet \code{SELECTION\_STATE} via Serial (COBS/CRC)
  \item Bridge empfängt, wandelt in JSON: \code{\{"type":"selection","selected":7\}}
  \item Browser empfängt via WebSocket, React aktualisiert UI
\end{enumerate}

% -----------------------------------------------------------------------------
\subsection{Repository-Struktur}
\label{subsec:repo-struktur}

Das Repository gliedert sich in drei Hauptbereiche: Firmware, Apps (Bridge und Dashboard)
sowie Hardware-Designs und Dokumentation.

\begin{lstlisting}[style=shell, caption={Repository-Struktur}, label={lst:repo-struktur}]
embedded_projekt/
|-- apps/
|   |-- bridge/                    # Python Bridge (Pi Runtime)
|   |   |-- bridge.py             # HTTP + WebSocket + Serial
|   |   |-- docs/
|   |   |   |-- PROTOCOL.md       # Serial/WS Protokoll (SSOT)
|   |   |   |-- WORKFLOW.md       # Build/Deploy/Run
|   |   |   `-- SMOKE_TEST.md     # E2E-Tests
|   |   `-- requirements.txt
|   |
|   `-- dashboard/                 # React Frontend
|       |-- src/
|       |   |-- pages/            # DashboardPage, SettingsPage
|       |   |-- components/       # UI-Komponenten
|       |   `-- lib/              # protocol.ts, ws.ts
|       |-- dist/                 # Build-Output (wird von Bridge served)
|       |-- docs/
|       |   |-- DEVELOPER.md      # Frontend-Entwicklung
|       |   `-- USER_GUIDE.md     # Benutzerhandbuch
|       |-- vite.config.ts
|       `-- package.json
|
|-- firmware/                      # ESP32-S3 Firmware (C++)
|   |-- include/hw/               # SSOT: Pins, Counts, SPI-Parameter
|   |-- src/
|   |   |-- hal/                  # SPI-Bus Zugriff
|   |   |-- drivers/              # CD4021, HC595
|   |   |-- logic/                # Debounce, Selection, Pipeline
|   |   |-- app/                  # Loop/Task Orchestrierung
|   |   `-- protocol/             # COBS/CRC Frames
|   |-- test/                     # Unit + E2E Tests
|   |-- docs/                     # Firmware-Dokumentation
|   `-- platformio.ini
|
|-- hardware/
|   `-- kicad/                    # Schaltplaene, PCB
|
|-- doc/
|   |-- md/                       # Markdown-Quellen (SSOT)
|   `-- latex/                    # LaTeX/PDF-Output
|
|-- assets/
|   |-- pdf/                      # Generierte PDFs
|   `-- video/                    # Medien
|
`-- scripts/                       # Build-/Deploy-Helfer
\end{lstlisting}

% -----------------------------------------------------------------------------
\subsection{SSOT -- Single Source of Truth}
\label{subsec:ssot}

\textbf{Regel:} Jede Konfiguration existiert an genau einer Stelle. Änderungen
propagieren automatisch.

\begin{table}[H]
\centering
\caption{SSOT-Dateien nach Aspekt}
\label{tab:ssot}
\begin{tabularx}{\textwidth}{@{}l X@{}}
\toprule
\textbf{Aspekt} & \textbf{SSOT-Datei} \\
\midrule
Pins                & \code{firmware/include/hw/board\_pins.h} \\
Counts/Masks        & \code{firmware/include/hw/io\_counts.h} \\
SPI-Parameter       & \code{firmware/include/hw/spi\_params.h} \\
Serial-Protokoll    & \code{apps/bridge/docs/PROTOCOL.md} \\
WebSocket-Protokoll & \code{apps/bridge/docs/PROTOCOL.md} \\
\bottomrule
\end{tabularx}
\end{table}

\begin{table}[H]
\centering
\caption{Generierte Artefakte}
\label{tab:artefakte}
\begin{tabularx}{\textwidth}{@{}l l X@{}}
\toprule
\textbf{Artefakt} & \textbf{Quelle} & \textbf{Output} \\
\midrule
Frontend & \code{apps/dashboard/src/} & \code{apps/dashboard/dist/} \\
PDFs     & \code{doc/latex/}          & \code{assets/pdf/} \\
\bottomrule
\end{tabularx}
\end{table}

% -----------------------------------------------------------------------------
\subsection{Firmware -- Build, Flash, Test}
\label{subsec:firmware}

\paragraph{Flashen}

\begin{lstlisting}[style=shell, caption={Firmware flashen}, label={lst:flash}]
cd /Users/jan/daten/IoT/embedded_projekt

# FreeRTOS-Variante (empfohlen)
pio run -d firmware -e seeed_xiao_esp32s3_freertos -t upload

# Arduino-Loop-Variante
pio run -d firmware -e seeed_xiao_esp32s3_loop -t upload
\end{lstlisting}

\paragraph{Monitor}

\begin{lstlisting}[style=shell]
pio device monitor -d firmware -e seeed_xiao_esp32s3_freertos
\end{lstlisting}

\paragraph{Tests}

\begin{lstlisting}[style=shell, caption={Unit- und E2E-Tests}, label={lst:tests}]
# Unit-Tests (Logik: Debounce, Selection, Pipeline)
pio test -d firmware -e seeed_xiao_esp32s3_test -v

# E2E-Hardwaretest (SPI, Pinning, Timing)
pio test -d firmware -e seeed_xiao_esp32s3_e2e -f test_e2e_hw -v
\end{lstlisting}

\begin{table}[H]
\centering
\caption{Build-Varianten}
\label{tab:build-varianten}
\begin{tabularx}{\textwidth}{@{}l X@{}}
\toprule
\textbf{Environment} & \textbf{Beschreibung} \\
\midrule
\code{seeed\_xiao\_esp32s3\_freertos} & FreeRTOS Task (\code{APP\_USE\_FREERTOS=1}) \\
\code{seeed\_xiao\_esp32s3\_loop}     & Arduino Loop (\code{APP\_USE\_FREERTOS=0}) \\
\code{seeed\_xiao\_esp32s3\_test}     & Unit-Tests (nur \code{logic/}) \\
\code{seeed\_xiao\_esp32s3\_e2e}      & E2E-Hardware-Test \\
\bottomrule
\end{tabularx}
\end{table}

% -----------------------------------------------------------------------------
\subsection{Dashboard -- Entwicklung und Build}
\label{subsec:dashboard}

\paragraph{DEV-Server (Mac)}

\begin{lstlisting}[style=shell]
cd /Users/jan/daten/IoT/embedded_projekt/apps/dashboard
npm ci
npm run dev
# -> http://localhost:5173/
# -> WebSocket via Proxy zu rover.local:8081
\end{lstlisting}

\paragraph{Production-Build}

\begin{lstlisting}[style=shell]
npm run lint
npm run typecheck
npm run build
# -> Output: dist/
\end{lstlisting}

\begin{table}[H]
\centering
\caption{npm Scripts}
\label{tab:npm-scripts}
\begin{tabularx}{\textwidth}{@{}l X@{}}
\toprule
\textbf{Script} & \textbf{Zweck} \\
\midrule
\code{npm run dev}       & DEV-Server mit HMR \\
\code{npm run lint}      & ESLint-Prüfung \\
\code{npm run typecheck} & TypeScript-Prüfung \\
\code{npm run build}     & Production-Build nach \code{dist/} \\
\code{npm run format}    & Prettier-Formatierung \\
\bottomrule
\end{tabularx}
\end{table}

% -----------------------------------------------------------------------------
\subsection{Bridge -- Raspberry Pi}
\label{subsec:bridge}

\paragraph{Manueller Start}

\begin{lstlisting}[style=shell, caption={Bridge manuell starten}, label={lst:bridge-start}]
ssh pi@rover
source ~/venvs/selection-bridge/bin/activate

python /home/pi/embedded_projekt/apps/bridge/bridge.py \
  --port auto \
  --vid 0x303A --pid 0x1001 \
  --baud 921600 \
  --http 8081 \
  --web-root /home/pi/embedded_projekt/apps/dashboard/dist \
  --wait-port --wait-timeout 0 --wait-poll 0.5 \
  --reconnect-delay 2 \
  --ack-timeout 0.6 --retries 4
\end{lstlisting}

\paragraph{systemd-Service}

\begin{lstlisting}[style=shell, caption={systemd-Befehle}, label={lst:systemd}]
# Status pruefen
sudo systemctl status embedded-bridge.service --no-pager -l

# Neu starten
sudo systemctl restart embedded-bridge.service

# Logs verfolgen
sudo journalctl -u embedded-bridge.service -f
\end{lstlisting}

\begin{table}[H]
\centering
\caption{URLs im PROD-Modus}
\label{tab:urls}
\begin{tabularx}{\textwidth}{@{}l X@{}}
\toprule
\textbf{Zweck} & \textbf{URL} \\
\midrule
Dashboard & \url{http://rover.local:8081/} \\
WebSocket & \url{ws://rover.local:8081/ws} \\
\bottomrule
\end{tabularx}
\end{table}

% -----------------------------------------------------------------------------
\subsection{Deployment-Workflow}
\label{subsec:deployment}

\paragraph{Mac → Pi synchronisieren}

\begin{lstlisting}[style=shell, caption={rsync-Synchronisation}, label={lst:rsync}]
# Zielpfad auf dem Pi anlegen (einmalig)
ssh pi@rover 'mkdir -p /home/pi/embedded_projekt/apps/bridge \
                       /home/pi/embedded_projekt/apps/dashboard'

# Sync (ohne venv, node_modules, dist)
rsync -avh --delete \
  --exclude '.git/' \
  --exclude '.venv/' \
  --exclude '__pycache__/' \
  --exclude '*.pyc' \
  --exclude '.DS_Store' \
  --exclude 'node_modules/' \
  --exclude 'apps/dashboard/dist/' \
  ./ \
  pi@rover:/home/pi/embedded_projekt/
\end{lstlisting}

\paragraph{Dashboard dist sync}

\begin{lstlisting}[style=shell, caption={Dashboard-Build und Sync}, label={lst:dist-sync}]
# Lokal bauen
cd apps/dashboard
npm ci
npm run build
cd -

# Dashboard dist sync (prod)
rsync -avh --delete \
  --exclude '.DS_Store' \
  ./apps/dashboard/dist/ \
  pi@rover:/home/pi/embedded_projekt/apps/dashboard/dist/
\end{lstlisting}

\paragraph{Service neu starten}

\begin{lstlisting}[style=shell]
ssh pi@rover "sudo systemctl daemon-reload && \
              sudo systemctl restart embedded-bridge.service"
\end{lstlisting}

% -----------------------------------------------------------------------------
\subsection{Dokumentationsstruktur}
\label{subsec:dokumentation}

\paragraph{Firmware-Dokumentation}

Die Firmware-Dokumentation liegt unter \code{firmware/docs/}:

\begin{table}[H]
\centering
\caption{Firmware-Dokumentation}
\label{tab:firmware-docs}
\begin{tabularx}{\textwidth}{@{}l X@{}}
\toprule
\textbf{Datei} & \textbf{Inhalt} \\
\midrule
\code{overview.md}     & Projektübersicht, Bit-Mapping (MSB-first) \\
\code{wiring.md}       & Hardware-Verdrahtung (Pins, ICs) \\
\code{architecture.md} & Schichtenmodell (HAL → Driver → Logic → App) \\
\bottomrule
\end{tabularx}
\end{table}

\paragraph{Bridge-Dokumentation}

Die Bridge-Dokumentation liegt unter \code{apps/bridge/docs/}:

\begin{table}[H]
\centering
\caption{Bridge-Dokumentation}
\label{tab:bridge-docs}
\begin{tabularx}{\textwidth}{@{}l X@{}}
\toprule
\textbf{Datei} & \textbf{Inhalt} \\
\midrule
\code{PROTOCOL.md}   & Serial + WebSocket Protokoll (SSOT) \\
\code{WORKFLOW.md}   & Build/Deploy/Run \\
\code{SMOKE\_TEST.md} & E2E-Tests \\
\bottomrule
\end{tabularx}
\end{table}

\paragraph{Dashboard-Dokumentation}

Die Dashboard-Dokumentation liegt unter \code{apps/dashboard/docs/}:

\begin{table}[H]
\centering
\caption{Dashboard-Dokumentation}
\label{tab:dashboard-docs}
\begin{tabularx}{\textwidth}{@{}l X@{}}
\toprule
\textbf{Datei} & \textbf{Inhalt} \\
\midrule
\code{DEVELOPER.md}  & Frontend-Entwicklung \\
\code{USER\_GUIDE.md} & Benutzerhandbuch \\
\bottomrule
\end{tabularx}
\end{table}

% -----------------------------------------------------------------------------
\subsection{Typische Arbeitsabläufe}
\label{subsec:workflows}

\begin{infobox}[Firmware-Logik ändern]
\begin{enumerate}
  \item Code ändern in \code{firmware/src/logic/}
  \item Unit-Tests: \code{pio test -d firmware -e seeed\_xiao\_esp32s3\_test -v}
  \item Bei Erfolg: Flashen mit \code{pio run ... -t upload}
  \item Monitor: \code{pio device monitor -d firmware}
\end{enumerate}
\end{infobox}

\begin{infobox}[Dashboard-UI ändern]
\begin{enumerate}
  \item DEV-Server starten: \code{cd apps/dashboard \&\& npm run dev}
  \item Ändern -- Browser aktualisiert automatisch (HMR)
  \item Lint + Build: \code{npm run lint \&\& npm run build}
  \item Deploy (siehe \cref{subsec:deployment})
\end{enumerate}
\end{infobox}

\begin{warnbox}[Hardware-Verdrahtung ändern]
Bei Änderungen an der Hardware-Verdrahtung:
\begin{enumerate}
  \item SSOT aktualisieren: \code{firmware/include/hw/board\_pins.h}
  \item Dokumentation aktualisieren: \code{firmware/docs/wiring.md}
  \item E2E-Test: \code{pio test ... -e seeed\_xiao\_esp32s3\_e2e -f test\_e2e\_hw -v}
\end{enumerate}
\end{warnbox}

% -----------------------------------------------------------------------------
\subsection{Roadmap}
\label{subsec:roadmap}

\begin{table}[H]
\centering
\caption{Projekt-Roadmap}
\label{tab:roadmap}
\begin{tabularx}{\textwidth}{@{}l c X@{}}
\toprule
\textbf{Milestone} & \textbf{Status} & \textbf{Beschreibung} \\
\midrule
10× Buttons/LEDs   & ✓  & Grundfunktion verifiziert \\
Serial-Protokoll   & ✓  & COBS/CRC, ACK/Retry \\
React-Dashboard    & ✓  & DEV/PROD-Split \\
100× Skalierung    & -- & Counts/Masks + Kaskadierung + Timing \\
Audio-Integration  & -- & Multimedia bei Selection-Change \\
\bottomrule
\end{tabularx}
\end{table}

% -----------------------------------------------------------------------------
\subsection{Siehe auch}
\label{subsec:projekt-siehe-auch}

Weiterführende Dokumentation:

\begin{itemize}
  \item \code{firmware/docs/overview.md} -- Firmware-Einstieg
  \item \code{apps/dashboard/docs/DEVELOPER.md} -- Frontend-Einstieg
  \item \code{apps/bridge/docs/WORKFLOW.md} -- Pi-Deployment
  \item \code{apps/bridge/docs/PROTOCOL.md} -- Protokoll-Spezifikation
\end{itemize}

% =============================================================================

\section{Selection-Panel-Projekt -- Gesamtübersicht}
\label{sec:projekt-uebersicht}

% -----------------------------------------------------------------------------
\subsection{Begriffsrahmen und Zielsetzung}
\label{subsec:begriffsrahmen}

\textbf{Projektroot:} \code{/Users/jan/daten/IoT/embedded\_projekt}

\textbf{Zielbild:} Ein interaktives Selection Panel mit 10--100 Buttons und LEDs,
bei dem zu jedem Zeitpunkt höchstens eine Auswahl aktiv ist (One-Hot-Prinzip).
Das System besteht aus drei Ebenen: Firmware (ESP32-S3), Bridge (Raspberry Pi),
Dashboard (Browser).

\begin{table}[H]
\centering
\caption{Zentrale Begriffe}
\label{tab:projekt-begriffe}
\begin{tabularx}{\textwidth}{@{}l X@{}}
\toprule
\textbf{Begriff} & \textbf{Definition} \\
\midrule
Selection Panel & Hardware-Interface mit Tastern und LEDs für Kanalauswahl \\
One-Hot         & Invariante: genau 0 oder 1 Auswahl aktiv, niemals mehrere \\
CD4021          & 8-Bit Parallel-zu-Seriell Schieberegister (Button-Input) \\
74HC595         & 8-Bit Seriell-zu-Parallel Schieberegister (LED-Output) \\
Bridge          & Python-Server, der WebSocket $\leftrightarrow$ Serial übersetzt \\
SSOT            & Single Source of Truth -- jede Konfiguration existiert an genau einer Stelle \\
\bottomrule
\end{tabularx}
\end{table}

% -----------------------------------------------------------------------------
\subsection{Systemarchitektur -- Drei Ebenen}
\label{subsec:systemarchitektur}

Die Architektur folgt einer Dreischichten-Struktur: Browser-Dashboard kommuniziert
via WebSocket mit der Bridge, die Bridge übersetzt zwischen JSON und binärem
Serial-Protokoll zur Firmware.

\begin{figure}[H]
\centering
\begin{tikzpicture}[
  layer/.style={draw, rectangle, rounded corners=3pt, minimum width=12cm,
                minimum height=2.0cm, font=\sffamily\small},
  arrow/.style={-{Stealth[length=3mm]}, thick},
  label/.style={font=\sffamily\footnotesize, text=MutedText}
]
  % Browser Layer
  \node[layer, fill=ArduinoTeal!10] (browser) at (0,6.5) {};
  \node[font=\sffamily\bfseries] at (-4.5,7.1) {BROWSER (Dashboard)};
  \node[font=\ttfamily\footnotesize] at (0,6.8) {React + TypeScript + Vite + Tailwind};
  \node[font=\ttfamily\footnotesize] at (0,6.2) {Quellcode: apps/dashboard/src/};

  % Bridge Layer
  \node[layer, fill=RaspberryRed!10] (bridge) at (0,3.25) {};
  \node[font=\sffamily\bfseries] at (-4.5,3.85) {RASPBERRY PI (Bridge)};
  \node[font=\ttfamily\footnotesize] at (0,3.55) {Python + aiohttp};
  \node[font=\ttfamily\footnotesize] at (0,3.0) {Runtime: apps/bridge/bridge.py};
  \node[font=\ttfamily\footnotesize] at (0,2.5) {Static: apps/dashboard/dist/};

  % Firmware Layer
  \node[layer, fill=SuccessGreen!10] (firmware) at (0,0) {};
  \node[font=\sffamily\bfseries] at (-4.5,0.6) {ESP32-S3 (Firmware)};
  \node[font=\ttfamily\footnotesize] at (0,0.3) {C++ + FreeRTOS};
  \node[font=\ttfamily\footnotesize] at (0,-0.3) {Hardware: CD4021 (Buttons) + 74HC595 (LEDs)};

  % Arrows with labels
  \draw[arrow] (browser.south) -- node[right, label] {WebSocket (JSON)}
               node[left, label] {ws://rover.local:8081/ws} (bridge.north);
  \draw[arrow] (bridge.south) -- node[right, label] {USB-CDC Serial (COBS + CRC16)}
               node[left, label] {\SI{921600}{Baud}} (firmware.north);
\end{tikzpicture}
\caption{Dreischichten-Architektur des Selection Panels}
\label{fig:architektur}
\end{figure}

\paragraph{Datenfluss bei Button-Druck}

Der Ablauf vom physischen Tastendruck bis zur UI-Aktualisierung durchläuft fünf Stufen:

\begin{enumerate}
  \item CD4021 liest Button-Matrix (alle \SI{10}{\milli\second})
  \item Firmware erkennt Press-Edge, aktualisiert \code{selectedId}
  \item Firmware sendet \code{SELECTION\_STATE} via Serial (COBS/CRC)
  \item Bridge empfängt, wandelt in JSON: \code{\{"type":"selection","selected":7\}}
  \item Browser empfängt via WebSocket, React aktualisiert UI
\end{enumerate}

% -----------------------------------------------------------------------------
\subsection{Repository-Struktur}
\label{subsec:repo-struktur}

Das Repository gliedert sich in drei Hauptbereiche: Firmware, Apps (Bridge und Dashboard)
sowie Hardware-Designs und Dokumentation.

\begin{lstlisting}[style=shell, caption={Repository-Struktur}, label={lst:repo-struktur}]
embedded_projekt/
|-- apps/
|   |-- bridge/                    # Python Bridge (Pi Runtime)
|   |   |-- bridge.py             # HTTP + WebSocket + Serial
|   |   |-- docs/
|   |   |   |-- PROTOCOL.md       # Serial/WS Protokoll (SSOT)
|   |   |   |-- WORKFLOW.md       # Build/Deploy/Run
|   |   |   `-- SMOKE_TEST.md     # E2E-Tests
|   |   `-- requirements.txt
|   |
|   `-- dashboard/                 # React Frontend
|       |-- src/
|       |   |-- pages/            # DashboardPage, SettingsPage
|       |   |-- components/       # UI-Komponenten
|       |   `-- lib/              # protocol.ts, ws.ts
|       |-- dist/                 # Build-Output (wird von Bridge served)
|       |-- docs/
|       |   |-- DEVELOPER.md      # Frontend-Entwicklung
|       |   `-- USER_GUIDE.md     # Benutzerhandbuch
|       |-- vite.config.ts
|       `-- package.json
|
|-- firmware/                      # ESP32-S3 Firmware (C++)
|   |-- include/hw/               # SSOT: Pins, Counts, SPI-Parameter
|   |-- src/
|   |   |-- hal/                  # SPI-Bus Zugriff
|   |   |-- drivers/              # CD4021, HC595
|   |   |-- logic/                # Debounce, Selection, Pipeline
|   |   |-- app/                  # Loop/Task Orchestrierung
|   |   `-- protocol/             # COBS/CRC Frames
|   |-- test/                     # Unit + E2E Tests
|   |-- docs/                     # Firmware-Dokumentation
|   `-- platformio.ini
|
|-- hardware/
|   `-- kicad/                    # Schaltplaene, PCB
|
|-- doc/
|   |-- md/                       # Markdown-Quellen (SSOT)
|   `-- latex/                    # LaTeX/PDF-Output
|
|-- assets/
|   |-- pdf/                      # Generierte PDFs
|   `-- video/                    # Medien
|
`-- scripts/                       # Build-/Deploy-Helfer
\end{lstlisting}

% -----------------------------------------------------------------------------
\subsection{SSOT -- Single Source of Truth}
\label{subsec:ssot}

\textbf{Regel:} Jede Konfiguration existiert an genau einer Stelle. Änderungen
propagieren automatisch.

\begin{table}[H]
\centering
\caption{SSOT-Dateien nach Aspekt}
\label{tab:ssot}
\begin{tabularx}{\textwidth}{@{}l X@{}}
\toprule
\textbf{Aspekt} & \textbf{SSOT-Datei} \\
\midrule
Pins                & \code{firmware/include/hw/board\_pins.h} \\
Counts/Masks        & \code{firmware/include/hw/io\_counts.h} \\
SPI-Parameter       & \code{firmware/include/hw/spi\_params.h} \\
Serial-Protokoll    & \code{apps/bridge/docs/PROTOCOL.md} \\
WebSocket-Protokoll & \code{apps/bridge/docs/PROTOCOL.md} \\
\bottomrule
\end{tabularx}
\end{table}

\begin{table}[H]
\centering
\caption{Generierte Artefakte}
\label{tab:artefakte}
\begin{tabularx}{\textwidth}{@{}l l X@{}}
\toprule
\textbf{Artefakt} & \textbf{Quelle} & \textbf{Output} \\
\midrule
Frontend & \code{apps/dashboard/src/} & \code{apps/dashboard/dist/} \\
PDFs     & \code{doc/latex/}          & \code{assets/pdf/} \\
\bottomrule
\end{tabularx}
\end{table}

% -----------------------------------------------------------------------------
\subsection{Firmware -- Build, Flash, Test}
\label{subsec:firmware}

\paragraph{Flashen}

\begin{lstlisting}[style=shell, caption={Firmware flashen}, label={lst:flash}]
cd /Users/jan/daten/IoT/embedded_projekt

# FreeRTOS-Variante (empfohlen)
pio run -d firmware -e seeed_xiao_esp32s3_freertos -t upload

# Arduino-Loop-Variante
pio run -d firmware -e seeed_xiao_esp32s3_loop -t upload
\end{lstlisting}

\paragraph{Monitor}

\begin{lstlisting}[style=shell]
pio device monitor -d firmware -e seeed_xiao_esp32s3_freertos
\end{lstlisting}

\paragraph{Tests}

\begin{lstlisting}[style=shell, caption={Unit- und E2E-Tests}, label={lst:tests}]
# Unit-Tests (Logik: Debounce, Selection, Pipeline)
pio test -d firmware -e seeed_xiao_esp32s3_test -v

# E2E-Hardwaretest (SPI, Pinning, Timing)
pio test -d firmware -e seeed_xiao_esp32s3_e2e -f test_e2e_hw -v
\end{lstlisting}

\begin{table}[H]
\centering
\caption{Build-Varianten}
\label{tab:build-varianten}
\begin{tabularx}{\textwidth}{@{}l X@{}}
\toprule
\textbf{Environment} & \textbf{Beschreibung} \\
\midrule
\code{seeed\_xiao\_esp32s3\_freertos} & FreeRTOS Task (\code{APP\_USE\_FREERTOS=1}) \\
\code{seeed\_xiao\_esp32s3\_loop}     & Arduino Loop (\code{APP\_USE\_FREERTOS=0}) \\
\code{seeed\_xiao\_esp32s3\_test}     & Unit-Tests (nur \code{logic/}) \\
\code{seeed\_xiao\_esp32s3\_e2e}      & E2E-Hardware-Test \\
\bottomrule
\end{tabularx}
\end{table}

% -----------------------------------------------------------------------------
\subsection{Dashboard -- Entwicklung und Build}
\label{subsec:dashboard}

\paragraph{DEV-Server (Mac)}

\begin{lstlisting}[style=shell]
cd /Users/jan/daten/IoT/embedded_projekt/apps/dashboard
npm ci
npm run dev
# -> http://localhost:5173/
# -> WebSocket via Proxy zu rover.local:8081
\end{lstlisting}

\paragraph{Production-Build}

\begin{lstlisting}[style=shell]
npm run lint
npm run typecheck
npm run build
# -> Output: dist/
\end{lstlisting}

\begin{table}[H]
\centering
\caption{npm Scripts}
\label{tab:npm-scripts}
\begin{tabularx}{\textwidth}{@{}l X@{}}
\toprule
\textbf{Script} & \textbf{Zweck} \\
\midrule
\code{npm run dev}       & DEV-Server mit HMR \\
\code{npm run lint}      & ESLint-Prüfung \\
\code{npm run typecheck} & TypeScript-Prüfung \\
\code{npm run build}     & Production-Build nach \code{dist/} \\
\code{npm run format}    & Prettier-Formatierung \\
\bottomrule
\end{tabularx}
\end{table}

% -----------------------------------------------------------------------------
\subsection{Bridge -- Raspberry Pi}
\label{subsec:bridge}

\paragraph{Manueller Start}

\begin{lstlisting}[style=shell, caption={Bridge manuell starten}, label={lst:bridge-start}]
ssh pi@rover
source ~/venvs/selection-bridge/bin/activate

python /home/pi/embedded_projekt/apps/bridge/bridge.py \
  --port auto \
  --vid 0x303A --pid 0x1001 \
  --baud 921600 \
  --http 8081 \
  --web-root /home/pi/embedded_projekt/apps/dashboard/dist \
  --wait-port --wait-timeout 0 --wait-poll 0.5 \
  --reconnect-delay 2 \
  --ack-timeout 0.6 --retries 4
\end{lstlisting}

\paragraph{systemd-Service}

\begin{lstlisting}[style=shell, caption={systemd-Befehle}, label={lst:systemd}]
# Status pruefen
sudo systemctl status embedded-bridge.service --no-pager -l

# Neu starten
sudo systemctl restart embedded-bridge.service

# Logs verfolgen
sudo journalctl -u embedded-bridge.service -f
\end{lstlisting}

\begin{table}[H]
\centering
\caption{URLs im PROD-Modus}
\label{tab:urls}
\begin{tabularx}{\textwidth}{@{}l X@{}}
\toprule
\textbf{Zweck} & \textbf{URL} \\
\midrule
Dashboard & \url{http://rover.local:8081/} \\
WebSocket & \url{ws://rover.local:8081/ws} \\
\bottomrule
\end{tabularx}
\end{table}

% -----------------------------------------------------------------------------
\subsection{Deployment-Workflow}
\label{subsec:deployment}

\paragraph{Mac → Pi synchronisieren}

\begin{lstlisting}[style=shell, caption={rsync-Synchronisation}, label={lst:rsync}]
# Zielpfad auf dem Pi anlegen (einmalig)
ssh pi@rover 'mkdir -p /home/pi/embedded_projekt/apps/bridge \
                       /home/pi/embedded_projekt/apps/dashboard'

# Sync (ohne venv, node_modules, dist)
rsync -avh --delete \
  --exclude '.git/' \
  --exclude '.venv/' \
  --exclude '__pycache__/' \
  --exclude '*.pyc' \
  --exclude '.DS_Store' \
  --exclude 'node_modules/' \
  --exclude 'apps/dashboard/dist/' \
  ./ \
  pi@rover:/home/pi/embedded_projekt/
\end{lstlisting}

\paragraph{Dashboard dist sync}

\begin{lstlisting}[style=shell, caption={Dashboard-Build und Sync}, label={lst:dist-sync}]
# Lokal bauen
cd apps/dashboard
npm ci
npm run build
cd -

# Dashboard dist sync (prod)
rsync -avh --delete \
  --exclude '.DS_Store' \
  ./apps/dashboard/dist/ \
  pi@rover:/home/pi/embedded_projekt/apps/dashboard/dist/
\end{lstlisting}

\paragraph{Service neu starten}

\begin{lstlisting}[style=shell]
ssh pi@rover "sudo systemctl daemon-reload && \
              sudo systemctl restart embedded-bridge.service"
\end{lstlisting}

% -----------------------------------------------------------------------------
\subsection{Dokumentationsstruktur}
\label{subsec:dokumentation}

\paragraph{Firmware-Dokumentation}

Die Firmware-Dokumentation liegt unter \code{firmware/docs/}:

\begin{table}[H]
\centering
\caption{Firmware-Dokumentation}
\label{tab:firmware-docs}
\begin{tabularx}{\textwidth}{@{}l X@{}}
\toprule
\textbf{Datei} & \textbf{Inhalt} \\
\midrule
\code{overview.md}     & Projektübersicht, Bit-Mapping (MSB-first) \\
\code{wiring.md}       & Hardware-Verdrahtung (Pins, ICs) \\
\code{architecture.md} & Schichtenmodell (HAL → Driver → Logic → App) \\
\bottomrule
\end{tabularx}
\end{table}

\paragraph{Bridge-Dokumentation}

Die Bridge-Dokumentation liegt unter \code{apps/bridge/docs/}:

\begin{table}[H]
\centering
\caption{Bridge-Dokumentation}
\label{tab:bridge-docs}
\begin{tabularx}{\textwidth}{@{}l X@{}}
\toprule
\textbf{Datei} & \textbf{Inhalt} \\
\midrule
\code{PROTOCOL.md}   & Serial + WebSocket Protokoll (SSOT) \\
\code{WORKFLOW.md}   & Build/Deploy/Run \\
\code{SMOKE\_TEST.md} & E2E-Tests \\
\bottomrule
\end{tabularx}
\end{table}

\paragraph{Dashboard-Dokumentation}

Die Dashboard-Dokumentation liegt unter \code{apps/dashboard/docs/}:

\begin{table}[H]
\centering
\caption{Dashboard-Dokumentation}
\label{tab:dashboard-docs}
\begin{tabularx}{\textwidth}{@{}l X@{}}
\toprule
\textbf{Datei} & \textbf{Inhalt} \\
\midrule
\code{DEVELOPER.md}  & Frontend-Entwicklung \\
\code{USER\_GUIDE.md} & Benutzerhandbuch \\
\bottomrule
\end{tabularx}
\end{table}

% -----------------------------------------------------------------------------
\subsection{Typische Arbeitsabläufe}
\label{subsec:workflows}

\begin{infobox}[Firmware-Logik ändern]
\begin{enumerate}
  \item Code ändern in \code{firmware/src/logic/}
  \item Unit-Tests: \code{pio test -d firmware -e seeed\_xiao\_esp32s3\_test -v}
  \item Bei Erfolg: Flashen mit \code{pio run ... -t upload}
  \item Monitor: \code{pio device monitor -d firmware}
\end{enumerate}
\end{infobox}

\begin{infobox}[Dashboard-UI ändern]
\begin{enumerate}
  \item DEV-Server starten: \code{cd apps/dashboard \&\& npm run dev}
  \item Ändern -- Browser aktualisiert automatisch (HMR)
  \item Lint + Build: \code{npm run lint \&\& npm run build}
  \item Deploy (siehe \cref{subsec:deployment})
\end{enumerate}
\end{infobox}

\begin{warnbox}[Hardware-Verdrahtung ändern]
Bei Änderungen an der Hardware-Verdrahtung:
\begin{enumerate}
  \item SSOT aktualisieren: \code{firmware/include/hw/board\_pins.h}
  \item Dokumentation aktualisieren: \code{firmware/docs/wiring.md}
  \item E2E-Test: \code{pio test ... -e seeed\_xiao\_esp32s3\_e2e -f test\_e2e\_hw -v}
\end{enumerate}
\end{warnbox}

% -----------------------------------------------------------------------------
\subsection{Roadmap}
\label{subsec:roadmap}

\begin{table}[H]
\centering
\caption{Projekt-Roadmap}
\label{tab:roadmap}
\begin{tabularx}{\textwidth}{@{}l c X@{}}
\toprule
\textbf{Milestone} & \textbf{Status} & \textbf{Beschreibung} \\
\midrule
10× Buttons/LEDs   & ✓  & Grundfunktion verifiziert \\
Serial-Protokoll   & ✓  & COBS/CRC, ACK/Retry \\
React-Dashboard    & ✓  & DEV/PROD-Split \\
100× Skalierung    & -- & Counts/Masks + Kaskadierung + Timing \\
Audio-Integration  & -- & Multimedia bei Selection-Change \\
\bottomrule
\end{tabularx}
\end{table}

% -----------------------------------------------------------------------------
\subsection{Siehe auch}
\label{subsec:projekt-siehe-auch}

Weiterführende Dokumentation:

\begin{itemize}
  \item \code{firmware/docs/overview.md} -- Firmware-Einstieg
  \item \code{apps/dashboard/docs/DEVELOPER.md} -- Frontend-Einstieg
  \item \code{apps/bridge/docs/WORKFLOW.md} -- Pi-Deployment
  \item \code{apps/bridge/docs/PROTOCOL.md} -- Protokoll-Spezifikation
\end{itemize}
